% !Mode:: "TeX:UTF-8"

\hitsetup{
  statesecrets={公开},
  natclassifiedindex={TP391},
  intclassifiedindex={68-05},
  %
  ctitleone={基于PyTorch-FSDP分布式训练},%本科生封面使用
  ctitletwo={的通信数据压缩研究},%本科生封面使用
  ctitlecover={基于PyTorch-FSDP分布式训练的通信数据压缩研究},%放在封面中使用，自由断行
  ctitle={基于PyTorch-FSDP分布式训练的通信数据压缩研究},%放在原创性声明中使用
  cxueke={工学},
  csubject={计算机科学与技术},
  caffil={深圳校区计算机科学与技术学院},
  cauthor={王云龙},
  csupervisor={施少怀},
  firstpagecdate={2026年5月},
  cdate={2026年5月},
  cstudentid={220110608},
  %
  etitle={Research on Communication Data Compression for PyTorch-FSDP Distributed Training},
  exueke={Engineering},
  esubject={Computer Science and Technology},
  eaffil={School of Computer Science and Technology, Shenzhen Campus},
  eauthor={Kuai Ge},
  esupervisor={Shi Shaohuai},
  edate={May, 2026},
  %
  ckeywords={分布式训练；全分片数据并行；通信压缩；梯度量化；梯度稀疏化；混合压缩；自适应调度},
  ekeywords={Distributed Training, Fully Sharded Data Parallel, Communication Compression, Gradient Quantization, Gradient Sparsification, Hybrid Compression, Adaptive Scheduling},
}

\begin{cabstract}

随着深度学习模型规模的持续增长，分布式训练已成为大规模模型训练的必要手段。PyTorch全分片数据并行（Fully Sharded Data Parallel, FSDP）通过将模型参数、梯度和优化器状态分片到多个设备上，有效降低了单设备的显存压力。然而，FSDP训练过程中节点间频繁的梯度同步通信成为制约训练效率的核心瓶颈。本文针对PyTorch FSDP框架下的通信优化问题，系统研究了通信数据压缩技术的实现方法、优化策略和性能影响。

首先，本文通过构建单卡与多卡对比实验，利用PyTorch Profiler对FSDP训练过程进行细粒度性能分析，定量识别了通信瓶颈：在GPT-2 Large模型双卡训练中，通信开销占单步耗时的61.2\%，且Reduce-Scatter操作位于反向传播关键路径上，阻塞效应显著。其次，本文设计了基于FSDP通信钩子机制的模块化通信压缩框架，将压缩策略与训练逻辑解耦，实现了量化类（int8、QSGD、Natural Compression、1-bit Seide）和稀疏化类（Top-k、Random-k、Threshold-v、Sketched-SGD）共9种压缩算法的统一集成，并在统一实验配置下对各算法进行了系统性对比评估。

在此基础上，本文进一步提出两项扩展研究。其一，设计并实现了两阶段混合压缩方法，将稀疏化与量化级联应用：第一阶段通过Top-k、Threshold-v或Random-k稀疏化筛选重要梯度分量，第二阶段对选中分量进行int8或1-bit量化，从而获得乘性压缩比。以Top-k（$\rho=0.01$）与int8级联为例，有效压缩比可达约44倍，远超单一方法。其二，提出自适应压缩调度策略，根据训练进度动态调整压缩强度：Warmup-Decay策略在训练初期使用较低压缩率保护梯度信息质量，随后逐步增大压缩率以加速通信；Linear策略按训练步数线性调整稀疏率；Grad-Adaptive策略通过跟踪梯度L2范数的指数移动平均并经Sigmoid映射，实现压缩强度随梯度统计特征的自动适配。

实验结果表明，在基线算法中，int8量化、Top-k和Threshold-v稀疏化在压缩比、通信时间和收敛质量的综合权衡上表现最为均衡：int8量化实现约4倍压缩比，训练吞吐提升约35\%，Loss仅上升3.8\%；Top-k稀疏化（1\%稀疏率）在保持良好收敛质量的同时达到最高吞吐。混合压缩方法在压缩比上实现了数量级提升，自适应调度策略在不额外牺牲收敛质量的前提下实现了更平稳的训练过程。本文的研究成果为构建高效、可扩展的大模型分布式训练系统提供了理论依据和工程参考。

\end{cabstract}

\begin{eabstract}

With the continuous growth of deep learning model scales, distributed training has become an essential approach for large-scale model training. PyTorch Fully Sharded Data Parallel (FSDP) effectively reduces single-device memory pressure by sharding model parameters, gradients, and optimizer states across multiple devices. However, the frequent gradient synchronization communication between nodes during FSDP training becomes a core bottleneck constraining training efficiency. This thesis systematically investigates the implementation methods, optimization strategies, and performance impacts of communication data compression techniques under the PyTorch FSDP framework.

First, this thesis quantitatively identifies the communication bottleneck through comparative experiments between single-GPU and multi-GPU configurations using PyTorch Profiler for fine-grained performance analysis of the FSDP training process. In GPT-2 Large model dual-GPU training, communication overhead accounts for 61.2\% of per-step time, with Reduce-Scatter operations on the critical path of backward propagation exhibiting significant blocking effects. Second, a modular communication compression framework based on the FSDP communication hook mechanism is designed, decoupling compression strategies from training logic, and achieving unified integration of 9 compression algorithms across quantization (int8, QSGD, Natural Compression, 1-bit Seide) and sparsification (Top-k, Random-k, Threshold-v, Sketched-SGD) categories. Systematic comparative evaluations of all algorithms are conducted under unified experimental configurations.

Building upon these foundations, this thesis further proposes two extensions. First, a two-stage hybrid compression method is designed and implemented, cascading sparsification with quantization: Stage-1 selects important gradient components via Top-k, Threshold-v, or Random-k sparsification, and Stage-2 applies int8 or 1-bit quantization to the selected components, achieving multiplicative compression ratios. Taking the cascade of Top-k ($\rho=0.01$) and int8 as an example, the effective compression ratio reaches approximately 44$\times$, far exceeding any single method. Second, adaptive compression scheduling strategies are proposed to dynamically adjust compression intensity according to training progress: the Warmup-Decay strategy uses lower compression rates in the early training phase to preserve gradient information quality, then gradually increases compression to accelerate communication; the Linear strategy adjusts the sparsity ratio linearly with training steps; the Grad-Adaptive strategy tracks the exponential moving average of gradient L2 norms and applies Sigmoid mapping, enabling automatic adaptation of compression intensity to gradient statistical characteristics.

Experimental results demonstrate that among baseline algorithms, int8 quantization, Top-k, and Threshold-v sparsification achieve the most balanced trade-off among compression ratio, communication time, and convergence quality: int8 quantization achieves approximately 4$\times$ compression ratio with 35\% throughput improvement and only 3.8\% Loss increase; Top-k sparsification (1\% sparsity ratio) achieves the highest throughput while maintaining good convergence quality. The hybrid compression methods achieve order-of-magnitude improvements in compression ratio, and the adaptive scheduling strategies enable smoother training processes without additional sacrifice in convergence quality. The research outcomes of this thesis provide theoretical foundations and engineering references for building efficient and scalable distributed training systems for large models.

\end{eabstract}
